\documentclass[../../../../main.tex]{subfiles}
\begin{document}

\begin{flushright}
\footnotesize\textit{【自動文字起こし・要確認】}
\end{flushright}

{\bf [解]}
\begin{align*}
\beta\gamma+\gamma\alpha+\alpha\beta=\frac12\left[(\alpha+\beta+\gamma)^2-(\alpha^2+\beta^2+\gamma^2)\right]=0
\end{align*}

だから，$\alpha\beta\gamma=k^3$とおくと，$\alpha,\beta,\gamma$は，$x$の3次方程式$x^3-k^3=0$の3解で，$\alpha\ne\beta$，$\gamma\ne\alpha$かつ$\alpha\ne0$である．
\begin{align*}
(x-k)(x^2+kx+k^2)=0
\end{align*}
\begin{align*}
\therefore\ x=k,\ \frac{-1\pm\sqrt3i}2k
\end{align*}

だから，$\triangle ABC$は正三角形（重心は原点）．

\newpage

\end{document}
